% Generated draft: requires mathematical and source review.
\documentclass[11pt]{article}
\usepackage[margin=25mm]{geometry}
\usepackage{fontspec}
\setmainfont{Latin Modern Roman}
\usepackage{amsmath,amssymb,mathtools,unicode-math}
\setmathfont{Latin Modern Math}
\usepackage{enumitem,needspace,xurl,hyperref}
\setlength{\parindent}{0pt}
\setlength{\parskip}{5pt}
\setlength{\emergencystretch}{3em}
\newcommand{\N}{\mathbb{N}}
\newcommand{\Z}{\mathbb{Z}}
\newcommand{\Q}{\mathbb{Q}}
\newcommand{\R}{\mathbb{R}}
\newcommand{\C}{\mathbb{C}}
\newcommand{\F}{\mathbb{F}}
\newcommand{\A}{\mathbb{A}}
\newcommand{\PP}{\mathbb{P}}
\newcommand{\E}{\mathbb{E}}
\newcommand{\eps}{\varepsilon}
\newcommand{\dd}{\,\mathrm{d}}
\begin{document}
\section*{UnsolvedMath ID: 20003093}
\textbf{Source problem code:} catalogue entry
\subsection{Definitions and mathematical statement}
Fix a finite-type orientable surface \(S\) in the hyperbolic range and its finite-area Teichm\"uller space. With Thurston's asymmetric metric
\[
d_{\rm Th}(X,Y)=\log\sup_{\gamma}\frac{\ell_\gamma(Y)}{\ell_\gamma(X)},
\]
define the geodesic envelope as the union of all directed geodesic segments from \(X\) to \(Y\). Equivalently, since geodesics exist, it is the metric interval
\[
\operatorname{Env}(X,Y)=\{Z:d_{\rm Th}(X,Z)+d_{\rm Th}(Z,Y)=d_{\rm Th}(X,Y)\}.
\]
The question attributed to K. Rafi asks whether this compact-set-valued map is continuous in both endpoints. Make the topology explicit by setting \(d_s(X,Y)=\max\{d_{\rm Th}(X,Y),d_{\rm Th}(Y,X)\}\) and using the Hausdorff metric of \(d_s\) on compact sets. In zero-dimensional Teichm\"uller space the assertion is immediate.

The supplied aid reports Hausdorff continuity for the once-punctured torus and for closed orientable surfaces of genus at least two. It says the literature checked did not cover all higher-complexity punctured surfaces in the original finite-type scope; this is a limitation of the supplied check, not a new assertion that no theorem exists. For the closed case it identifies Pan--Wolf's headline Theorem 1.1 as unrestricted and detailed Theorem 5.1 as stated off the diagonal.

The aid supplies an abstract diagonal argument. Let \((M,d)\) be a directed geodesic asymmetric metric space, with continuous distance, compact closed forward balls in the symmetrized topology, and compact-set reversibility: for every compact \(K\) there is \(C_K\ge1\) such that
\[
d(b,a)\le C_Kd(a,b)\qquad(a,b\in K).
\]
Let \(I_d(x,y)=\{z:d(x,z)+d(z,y)=d(x,y)\}\). Its intervals are compact and upper semicontinuous in endpoints: if \(x_n\to x,y_n\to y\), no points of \(I_d(x_n,y_n)\) can remain a positive symmetrized distance from \(I_d(x,y)\). For \(x_n,y_n\to x\), a single suitable compact \(K\) contains the relevant intervals eventually, and
\[
d_H^s\bigl(I_d(x_n,y_n),\{x\}\bigr)
\le d_s(x,x_n)+C_Kd(x_n,y_n)\longrightarrow0.
\]
Thus the aid obtains diagonal continuity for Thurston envelopes under these standard finite-type metric hypotheses, quantitatively separating it from the off-diagonal issue. Upper semicontinuity alone does not imply full continuity: lower semicontinuity requires each point of the limiting interval to be approximated by points of nearby intervals. For punctured cases not covered by the supplied references, that off-diagonal lower-semicontinuity step remains the unresolved target in the report.
\subsection{Short English statement}
Ask whether all Thurston geodesics between two endpoints form a Hausdorff-continuous envelope. The supplied aid reports continuity in stated surface classes and proves diagonal collapse under properness and local reversibility, while distinguishing this from off-diagonal lower semicontinuity in unverified punctured cases.
\subsection{Sources}
\begin{itemize}
\item K. Rafi, Problem 2.3, AIM workshop \emph{Lipschitz metric on Teichmueller space}, \url{https://aimath.org/pastworkshops/lipschitzteichproblems.pdf}.
\item The research aid cites Pan--Wolf, headline Theorem 1.1 and detailed Theorem 5.1, for closed surfaces, and refers to known once-punctured-torus continuity without a full bibliographic entry for either source.
\item Local UnsolvedMath registry, permanent ID 20003093, exact code \texttt{catalogue entry}; canonical location \texttt{aim-topology-notes.json}, notes[180], dated 14 August 2026; statement labeled exact.
\item Research attempt 1: partial result, problem partially solved at the run, rigor label proved, novelty confidence low; registry status partially solved. The abstract lemma and literature-scope account are machine-generated research aids, not peer reviewed; independent verification is required before citation.
\end{itemize}
\end{document}
